\cgeentry{cge_multiplier}{SAM multipliers and their decomposition}

\cgesec{Description}

\cmd{cge_multiplier} computes the accounting multipliers of the SAM mapped by
\manref{cge_map} (Pyatt and Round 1979): the effects of an injection into one
account on the incomes of all the endogenous accounts, the accounts whose
spending the SAM explains, with prices fixed. For each injected account it
displays the total multiplier and its parts in gross output, value added and
household income. \opt{decomp}{ose} splits the multipliers into transfer,
open-loop and closed-loop effects, and \opt{inj}{ect()} gives the effects of
a vector of injections.

The multipliers use the SAM after the layout rewrites of \manref{cge_map} and
need no calibration. The model in memory is not changed.

\cgesec{Quick start}

Multipliers of every endogenous account
\begin{stcode}
cge_multiplier
\end{stcode}
As above, with the decomposition into transfer, open-loop and closed-loop
effects
\begin{stcode}
cge_multiplier, decompose
\end{stcode}
Effects of 10 more spending on health services
\begin{stcode}
cge_multiplier, inject(com_hlt:10)
\end{stcode}
Multipliers with the government endogenous, loaded as data
\begin{stcode}
cge_multiplier, endogenous(gov) clear
\end{stcode}

\cgesec{Syntax}

\begin{stsyntax}
\cmdb{cge_multiplier} [\texttt{,} \argu{options}]
\end{stsyntax}

\begin{opttable}
\optgroup{Accounts}
\optarg{endo}{genous}{accounts} & make exogenous accounts endogenous, in the institutions block\\
\optarg{exo}{genous}{accounts} & make endogenous accounts exogenous\\
\optgroup{Results}
\opt{decomp}{ose} & add the decomposition into transfer, open-loop and closed-loop effects\\
\optarg{inj}{ect}{\argu{account}\texttt{:}\argu{value} \dots} & display the effects of injections into endogenous accounts\\
\optarg{acc}{ounts}{accounts} & injected accounts to display\\
\opt{}{clear} & replace the data in memory with the multipliers or the effects\\
\end{opttable}

\cgesec{Options}

\begin{optdesc}
\item[\optarg{endo}{genous}{accounts}] adds accounts that are exogenous by
default, for example the government, to the endogenous accounts, in the
institutions block.

\item[\optarg{exo}{genous}{accounts}] makes accounts that are endogenous by
default exogenous, for example an enterprise.

\item[\opt{decomp}{ose}] adds the decomposition of each multiplier into
transfer, open-loop and closed-loop effects. It requires the SAM to have no
flows among the blocks other than those of the circular flow; the command lists
the flows that break the structure if it does.

\item[\optarg{inj}{ect}{\argu{account}\texttt{:}\argu{value} \dots}] gives
injections into endogenous accounts, separated by spaces, for example
\cmd{inject(com_hlt:100)} for 100 more spending on health services by an
exogenous account such as the government. The command then displays, instead
of the multipliers, the change in each endogenous account's income, its income
in the SAM and the percentage change, and the totals.

\item[\optarg{acc}{ounts}{accounts}] chooses the injected accounts whose
multipliers are displayed. The default is every endogenous account, or the
commodities and households if there are more than 40.

\item[\opt{}{clear}] replaces the data in memory with the multiplier matrix,
one observation per pair of endogenous accounts, with the variables
\cmd{row}, \cmd{column}, \cmd{row_block}, \cmd{column_block} and
\cmd{multiplier}, and, with \opt{decomp}{ose}, \cmd{transfer},
\cmd{open_loop} and \cmd{closed_loop}. With \opt{inj}{ect()}, it loads instead the effects of
the injection, one observation per endogenous account, with the variables
\cmd{account}, \cmd{block}, \cmd{base}, \cmd{change} and \cmd{pct}.
\end{optdesc}

\cgesec{Remarks and examples}

The accounts of the SAM are divided into endogenous accounts, whose incomes
are determined by the spending of the SAM's accounts, and exogenous accounts,
whose spending injects income into them. By default the endogenous accounts
form three blocks:
\begin{cgetable}{lp{0.70\linewidth}}
\toprule
Block & Accounts\\
\midrule\endhead
\bottomrule\endlastfoot
production & activities, commodities and margin accounts\\
factors & factors\\
institutions & households and enterprises, and the accounts in \opt{endo}{genous()}\\
\end{cgetable}
The government, savings--investment, the rest of the world, the tax accounts
and the health insurance fund are exogenous. Each endogenous account spends
its income in the shares of its column of the SAM; spending on exogenous
accounts (taxes, savings and imports) leaks out of the circular flow.

Accounting multipliers assume fixed prices, unused capacity in every sector,
and average propensities to spend equal to marginal ones. They are upper
bounds on what a CGE model with supply constraints and price changes would
find, and complement it.

\cgesubsec{Example 1: the multipliers}

The health layout of the synthetic SAM is mapped:
\stlog{mult_setup}
\cgeneed{18\baselineskip}
The multipliers have one row per injected account:
\stlog[\scriptsize]{mult_table}
A unit injection into health services, \cmd{com_hlt}, raises the incomes of
all endogenous accounts by 9.408 in total, including the injection itself:
gross output by 2.581, value added by 1.864 and household income by 1.813.
Health services are neither imported nor taxed, so the whole of the injection
passes to the health activity, and the multipliers of \cmd{com_hlt} are those
of \cmd{act_hlt}, plus the injection itself in the total.

\cgesubsec{Example 2: the decomposition}

With \opt{decomp}{ose}, each total multiplier is split into 1, the injection
itself, and three effects:
\stlog[\scriptsize]{mult_decomp}
The transfer effects are those within the injected account's block, such as
the intermediate purchases among activities and commodities; the open-loop
effects are those on the other blocks, before any income returns to the
injected block; and the closed-loop effects are those of the circular flow
production $\to$ factors $\to$ institutions $\to$ production back to it. The
factors and the household have no transfer effects: no factor pays another,
and the household pays no other institution.

\cgesubsec{Example 3: an injection}

An injection of 10 into health services, such as more spending on them by the
government or another exogenous account:
\stlog{mult_inject}
Gross output rises by 25.8 and household income by 18.1. The base income of
an account is its column total in the SAM.

\cgesubsec{Example 4: other endogenous accounts}

\opt{endo}{genous()} and \opt{exo}{genous()} change the accounts whose incomes
the multipliers determine. With the government endogenous:
\stlog[\scriptsize]{mult_endo}
The multipliers of \cmd{com_hlt} and \cmd{hhd} do not change: the government's
income comes only from the tax accounts, which remain exogenous, so no
injection elsewhere reaches it. Making the government endogenous adds the
effects of an injection into its income, through its purchases and its
transfers to the household.

Accounts added to the institutions block may break the circular structure
that the decomposition requires. A sales tax account paid by the commodities
is a flow from production to institutions:
\stlog{mult_offcircle}

\cgesubsec{Example 5: the multipliers as data}

With \opt{}{clear}, the matrix of multipliers replaces the data in memory. The
column of \cmd{com_hlt} gives the effects of a unit injection into it on each
account:
\stlog{mult_clear}
Off the diagonal, each multiplier is the sum of its transfer, open-loop and
closed-loop parts; on the diagonal, the sum is one less than the multiplier.

\cgesec{Stored results}

\cmd{cge_multiplier} stores the following in \cmd{r()}:
\begin{storedres}
\resgroup{Scalars}
\res{r(N_endogenous)} & number of endogenous accounts\\
\res{r(injection)} & total injection (with \opt{inj}{ect()})\\
\res{r(output)} & change in gross output (with \opt{inj}{ect()})\\
\res{r(value_added)} & change in value added (with \opt{inj}{ect()})\\
\res{r(household_income)} & change in household income (with \opt{inj}{ect()})\\
\res{r(total)} & change in the incomes of all endogenous accounts (with \opt{inj}{ect()})\\
\resgroup{Macros}
\res{r(endogenous)} & endogenous accounts\\
\end{storedres}

\cgesec{Methods and formulas}

Let $N$ be the set of endogenous accounts and $t_j$ the column total of
account $j$ in the SAM after the layout rewrites. The average expenditure
propensities among the endogenous accounts are
\begin{equation*}
A_{ij} = \frac{\SAM_{ij}}{t_j}, \qquad i, j \in N,
\end{equation*}
with $A_{ij} = 0$ when $t_j = 0$. The incomes $y$ of the endogenous accounts
satisfy $y = A y + x$, where $x$ is their income from the exogenous accounts,
so that
\begin{equation*}
y = M x, \qquad M = (I - A)^{-1}.
\end{equation*}
Column $j$ of $M$ is the change in the income of every endogenous account for
a unit injection into account $j$. $I - A$ is invertible when every chain of
payments among the endogenous accounts eventually leaks to an exogenous
account; if it is singular, \cmd{cge_multiplier} stops and asks for one of the
accounts to be made exogenous. For injected account $j$ the command displays
the sums of column $j$ of $M$ over all endogenous accounts (\cmd{Total}), over
the activities (\cmd{Output}), over the factors (\cmd{Value add.}) and over the
households (\cmd{Hh income}). With \opt{inj}{ect()}, the changes
in income are $\Delta y = M x$ for the vector $x$ of injections, and the
percentage change of account $i$ is $100\, \Delta y_i / |t_i|$.

\textit{Decomposition.} With the blocks ordered as production (1), factors (2)
and institutions (3), the circular flow is made of the payments of production
to factors ($A_{21}$), of factors to institutions ($A_{32}$) and of
institutions to production ($A_{13}$), with transfers within each block
($A_{11}$, $A_{22}$, $A_{33}$). The decomposition requires the other blocks,
$A_{12}$, $A_{23}$ and $A_{31}$, to be zero. With $\tilde A$ the
block-diagonal part of $A$ and $A^* = (I - \tilde A)^{-1}(A - \tilde A)$,
\begin{equation*}
M = M_3 M_2 M_1, \qquad M_1 = (I - \tilde A)^{-1}, \quad
M_2 = I + A^* + A^{*2}, \quad M_3 = \bigl(I - A^{*3}\bigr)^{-1},
\end{equation*}
the transfer, open-loop and closed-loop multipliers of Pyatt and Round (1979),
and additively
\begin{equation*}
M = I + (M_1 - I) + (M_2 - I) M_1 + (M_3 - I) M_2 M_1 .
\end{equation*}
The columns \cmd{Transfer}, \cmd{Open loop} and \cmd{Closed} are the column
sums of the last three terms, so that the total multiplier is one plus their
sum. Multipliers and their decomposition are described in
\secref{sec:multipliers} and \secref{sec:decomposition}.

\cgesec{References}

\begin{list}{}{\setlength{\leftmargin}{2em}\setlength{\itemindent}{-2em}%
  \setlength{\itemsep}{2pt}\setlength{\parsep}{0pt}}
\item Pyatt, G., and J. I. Round. 1979. Accounting and fixed price multipliers
in a social accounting matrix framework. \textit{Economic Journal} 89:
850--873.
\end{list}

\cgesec{Also see}

\begin{alsosee}
\seeentry{cge_spa}{Structural path analysis of a SAM multiplier}
\seeentry{cge_map}{Assign SAM accounts to model roles}
\seeentry{cge_sam}{Read a social accounting matrix}
\seeentry{cge}{Computable general equilibrium modelling in Stata}
\seechapter{chap:samtools}{SAM tools}
\end{alsosee}
